\section{Throughput Analysis}
\label{sec:analysis}

\subsection{Per-Vertex Processing Time}

The total cycles to process vertex $v$ with degree $d_v$ is:
\begin{equation}
T(v) = T_\text{dequeue} + T_\text{ptr} + T_\text{edge}(d_v) + T_\text{scatter}(d_v)
\label{eq:Tv}
\end{equation}

where the components are:

\begin{itemize}[leftmargin=*,itemsep=2pt]
\item $T_\text{dequeue} \approx 2$ cycles: FIFO read handshake.
\item $T_\text{ptr} = T_\text{access} + T_\text{fsmov} \approx 20 + 3 = 23$ cycles:
  AXI4 round trip for row pointer fetch (1 beat) plus FSM state transitions
  (\textsc{FetchPtr} $\to$ \textsc{WaitPtr} $\to$ \textsc{CheckEdges}).
\item $T_\text{edge}(d_v) = \lceil d_v / 8 \rceil \cdot T_\text{access}$: AXI4 burst
  for adjacency list. Each beat delivers 8 neighbor IDs; the burst issues once and
  streams $\lceil d_v/8 \rceil$ beats back.
\item $T_\text{scatter}(d_v) = d_v \cdot T_\text{RMW}$: bitmap check + optional set
  per neighbor. $T_\text{RMW} = 2$ cycles (1 BRAM read pipeline stage + 1 evaluate).
\end{itemize}

\subsection{Steady-State Throughput}

For graphs where the average degree $\bar{d}$ is large, the edge burst amortizes
the AXI latency across many neighbors:

\begin{equation}
T_\text{edge}(\bar{d}) = \left\lceil \frac{\bar{d}}{8} \right\rceil \cdot T_\text{access}
\approx \frac{\bar{d}}{8} \cdot T_\text{access}
\quad \text{(for } \bar{d} \gg 8\text{)}
\end{equation}

Cycles per edge in steady state:
\begin{equation}
\frac{T(v)}{d_v} \approx
\frac{T_\text{ptr}}{d_v}
+ \frac{T_\text{access}}{8}
+ T_\text{RMW}
\label{eq:cpe}
\end{equation}

For $\bar{d} \gg T_\text{ptr}$ (high-degree hubs), the pointer fetch overhead
vanishes and cycles-per-edge approaches:

\begin{equation}
\frac{T(v)}{d_v} \bigg|_{\bar{d} \to \infty}
= \frac{T_\text{access}}{8} + T_\text{RMW}
= \frac{20}{8} + 2 = 4.5 \text{ cycles/edge}
\label{eq:cpe_limit}
\end{equation}

For the opposite extreme, low-degree vertices ($d_v = 1$):
\begin{equation}
T(v)\big|_{d_v=1} \approx 23 + 20 + 2 = 45 \text{ cycles/vertex}
\label{eq:low_degree}
\end{equation}

\subsection{Connection to HBM Bandwidth Efficiency}

The companion analysis~\cite{fritz2026rmwwall} shows that graph BFS achieves
$\eta \approx 40\%$ HBM bandwidth efficiency on HBM4, dominated by the
$t_\text{RRD,S}$ row-activation scheduling bottleneck with row-buffer hit rate
$\lambda = 0.06$ for power-law graphs.

This result has a direct microarchitectural interpretation in BFS-HBM.
Consider the fraction of cycles spent on useful data transfer vs.\ total cycles:

\begin{equation}
\eta_\text{arch} = \frac{T_\text{RMW} \cdot \bar{d}}{T(v)}
\approx \frac{2\bar{d}}{23 + \frac{T_\text{access}\bar{d}}{8} + 2\bar{d}}
\label{eq:eta_arch}
\end{equation}

For $\bar{d} = 10$ (a typical web-scale power-law graph~\cite{leskovec2009community}):

\begin{equation}
\eta_\text{arch}(10) = \frac{20}{23 + 25 + 20} = \frac{20}{68} \approx 29\%
\end{equation}

The gap between $\eta_\text{arch} = 29\%$ and the system-level $\eta = 40\%$ reflects
two effects operating in the same direction: the JEDEC timing model in~\cite{fritz2026rmwwall}
characterizes \emph{bus utilization efficiency} (useful data bytes / total bytes transferred),
while $\eta_\text{arch}$ measures \emph{compute utilization} (cycles doing RMW /
total cycles). Both are limited by the irregular access pattern; the precise values
depend on graph topology, HBM channel count, and operating clock.

\subsection{The Single-Outstanding-Read Bottleneck}

The most significant architectural constraint is that BFS-HBM issues one AXI4 read
transaction at a time. The FSM cannot issue the adjacency list burst until the row
pointer read completes, and it cannot begin the next vertex's row pointer read
until the scatter loop for the current vertex completes.

This serialization is not an accidental artifact of the current FSM design — it is
a consequence of the data dependency chain in BFS:

\begin{equation}
\mathit{scatter}(v) \to \mathit{enqueue}(\text{neighbors}) \to \mathit{dequeue}(u)
\to \mathit{fetch\_ptr}(u) \to \cdots
\end{equation}

The hop-$k$ frontier is logically unavailable until hop-$k-1$ is fully expanded.
This is the same causal constraint that limits pipeline depth in any sequentially
dependent computation.

Breaking this constraint requires either:
\begin{itemize}[leftmargin=*,itemsep=2pt]
\item \textbf{Multiple parallel traversal engines}: $P$ independent
  \texttt{bfs\_ctrl} instances sharing the HBM arbiter, each processing a
  different frontier vertex simultaneously. Throughput scales as $P$ for
  non-overlapping frontiers; an arbiter prevents channel conflicts.
\item \textbf{Direction-optimizing BFS}~\cite{beamer2012direction}: switch from
  top-down (expand frontier) to bottom-up (check all unvisited vertices against frontier)
  when the frontier becomes large. Reduces pointer-chasing for high-diameter graphs.
\item \textbf{Prefetch buffering}: issue the next vertex's row pointer read
  concurrently with the scatter loop of the current vertex, using a 1-deep
  prefetch register. This requires tracking the next frontier head speculatively
  and is left as an extension.
\end{itemize}

Section~\ref{sec:silicon} discusses multi-engine parallelism as the primary
throughput scaling mechanism for the ASIC design.
